\documentclass[10pt]{article}
\usepackage[margin=20mm]{geometry}
\usepackage{amsmath,amssymb,mathtools,fontspec,unicode-math,xurl,hyperref}
\IfFontExistsTF{DejaVu Serif}{\setmainfont{DejaVu Serif}}{\setmainfont{Latin Modern Roman}}
\setmathfont{latinmodern-math.otf}
\setlength{\parindent}{0pt}\setlength{\parskip}{6pt}\sloppy
\title{Ghosts of abelian schemes: pointwise descent, Fourier purity, and auxiliary hypotheses}
\author{AI-assisted technical review}
\date{}
\begin{document}
\maketitle
Abstract: Four findings are established in the supplied author manuscript of Cadoret and Tamagawa. A quadratic Kummer system disproves the pointwise submodule conclusion of Proposition 5 when only a constant-field extension is allowed. An elementary finite-monodromy sheaf with a deleted stalk disproves the unrestricted Fourier-realization theorem in Appendix A. Two further local issues concern the hypotheses of a density lemma and the claimed uniqueness of conjugation between maximal tori. The examples do not refute the central maximality statement formulated in the almost category, or the geometric isogeny applications after appropriate corrections.\par
\section*{Identity, conventions and review scope}
The reviewed file is Anna Cadoret and Akio Tamagawa, Ghosts and families of abelian varieties with a common isogeny factor, 26-page author-hosted manuscript [main]. Its revision date is undetermined. The journal identity is J. Eur. Math. Soc., DOI 10.4171/JEMS/1597 (https://doi.org/10.4171/JEMS/1597), online first 27 January 2025. Equivalence of the author file with the publisher article has not been established. Page references below are one-based PDF pages, equal to the manuscript’s printed pages.\par
The paper explicitly distinguishes actual representations from the almost category obtained by restricting to open subgroups (p. 2). It also distinguishes arithmetic monodromy \(G_\ell\) from geometric monodromy \(\overline G_\ell\). These distinctions are essential to the first finding. The review follows the definitions, Hilbertian argument, density lemmas, ghost construction and gcd argument, abelian realization and Appendix. Detailed arithmetic checks are separated below from imported Tate, Honda–Tate, monodromy, companion, and weight theorems.\par
\section*{Constant-field extension does not give every pointwise factor}
Problem 1 (Major: the last assertion of Proposition 5 is too strong).  Proposition 5, p. 4, allows a finite extension of \(k\) and concludes that, for every closed point \(s\), \(a_s^*\Psi(V_\ell)\) is an actual submodule of \((s^*V_\ell)^{\mathrm{ss}}\). This conclusion fails for a rank-one pure rational compatible system.\par
Take \(k=\mathbb F_q\) with \(q\) odd, \(S=\mathbb G_m/k\), and the connected finite étale cover \[\mathbb G_m\longrightarrow\mathbb G_m,\qquad u\longmapsto u^2.\] For every \(\ell\ne p\), let \(V_\ell\) be its sign representation over \(\mathbb Q_\ell\). It has rank one, finite monodromy, rational characteristic polynomials \(T-1\) or \(T+1\) independent of \(\ell\), and weight zero. Its geometric algebraic monodromy is \(\mu_2\), whose maximal torus is trivial. Consequently \[\dim V_\ell^{\overline T_\ell}=1.\]\par
After any constant extension \(K=\mathbb F_{q^m}\), choose a square \(a\in K^\times\) and a nonsquare \(b\in K^\times\). They give two \(K\)-rational closed points of \(S_K\). The Frobenius eigenvalues of the sign sheaf at these points are respectively \(+1\) and \(-1\); inversion from arithmetic to geometric Frobenius makes no difference. If a one-dimensional \(\pi_1(K)\)-module \(W\) were an actual submodule at both points, then, since both points have degree one and both local representations have rank one, \[F_K|W=1\qquad\text{and}\qquad F_K|W=-1.\] This is impossible. The counterexample remains valid for every allowed constant-field extension, and satisfies all the purity and compatibility assumptions of the proposition.\par
\subsection*{Location of the scope change and repair}
Section 6.2.1, p. 12, passes to an arbitrary connected étale cover of \(S\) to obtain connected monodromy and a torsion-free group generated by Frobenius eigenvalues. Section 6.3.2, p. 13, uses those stronger hypotheses to put every semisimple Frobenius in a maximal torus and compare its action. Those reductions are appropriate for almost statements, but do not descend the asserted actual embedding to the original \(S\) after merely extending its constants.\par
There are two distinct possible corrections: retain the actual embedding only on a suitable finite étale cover of \(S\), with the reduction hypotheses explicitly incorporated; or formulate the original-base comparison as an almost submodule statement, with the needed descent argument. In this example the double cover trivializes the system, and at any fixed closed point the eigenvalue \(-1\) becomes \(+1\) after a quadratic residue-field extension. Thus the almost comparison survives.\par
The part of Proposition 5 quantified over conjugating elements \(g\) is not separately contradicted here: a nonidentity element of \(\mu_2\) cannot be conjugated into the trivial torus, so that conditional quantifier is empty at the relevant point. The decisive contradiction concerns the unconditional final assertion. This is a Major scope error in a central proposition, not a disproof of the existence of the ghost as an almost representation, Theorem 6’s almost maximality, or every geometric isogeny conclusion.\par
\section*{An elementary Fourier transform with a weight-zero summand}
Problem 2 (Major: Theorem 22 lacks a middle-extension or purity condition).  Appendix A, p. 23, assumes only that \(\mathcal F\) is elementary in the stated sense and has finite monodromy on a nonempty open set. It concludes that, on some nonempty open \(U_2\), its naive Fourier transform is a direct factor of \(R^1\pi_*\overline{\mathbb Q}_\ell\) for an abelian scheme. The conclusion is false even when the generic rank of \(\mathcal F\) is one.\par
Take \(k=\mathbb F_3\), \(\ell=2\), a nontrivial additive character \(\psi\) with values in \(\overline{\mathbb Q}_2^\times\), and \[j:\mathbb G_m\hookrightarrow\mathbb A^1,\qquad i:\{0\}\hookrightarrow\mathbb A^1,
 \qquad L=\mathcal L_{\psi(t^2)},\qquad \mathcal F=j_!j^*L.\] The generic monodromy has finite order three and is trivialized by an Artin–Schreier cover. The sheaf has no nonzero punctual sections, so it satisfies (Elem-1).\par
For every \(a\in\overline k\), \[H_c^2(\mathbb A^1_{\overline k},\mathcal F\otimes\mathcal L_{\psi(at)})
   =H_c^2({\mathbb G_m}_{\overline k},\mathcal L_{\psi(t^2+at)})=0.\] Indeed, the polynomial \(t^2+at\) has at infinity a pole of order two, prime to \(p=3\), and defines a geometrically nontrivial rank-one Artin–Schreier sheaf. A rational expression \(h^3-h\) with a pole has pole order divisible by three, so it cannot eliminate that pole. Poincaré duality identifies the indicated top compactly supported cohomology with the dual invariant space, which vanishes. Thus (Elem-2) holds for every \(a\), as required.\par
There is a short exact sequence \[0\longrightarrow\mathcal F\longrightarrow L
       \longrightarrow i_*\overline{\mathbb Q}_2\longrightarrow0.\] Tensor on \(\mathbb A^1_t\times\mathbb A^1_x\) with the Fourier kernel \(\mathcal L_{\psi(tx)}\) and apply \(Rp_{2!}\). The degree-zero compactly supported cohomology of the smooth rank-one sheaf \(\mathcal L_{\psi(t^2+xt)}\) is zero. The transform in degree zero of the punctual quotient is the constant sheaf. The long exact sequence therefore gives \[0\longrightarrow\overline{\mathbb Q}_2
    \longrightarrow \mathcal F_\psi(\mathcal F)
    \longrightarrow \mathcal F_\psi(L)\longrightarrow0.\] In particular, the Fourier transform contains a nonzero constant sub-local system on every nonempty open where it is lisse. At every closed point this contributes Frobenius eigenvalue \(1\).\par
For an abelian scheme \(\pi:A\to U_2\), the eigenvalues of geometric Frobenius on \((R^1\pi_*\overline{\mathbb Q}_2)_s\) have complex absolute value \(|k(s)|^{1/2}>1\). It cannot contain that constant subrepresentation, and hence cannot contain the proposed Fourier transform even as a subobject, much less as a direct factor. Shrinking \(U_2\) cannot remove a constant sub-local system.\par
\subsection*{The cited purity theorem and the repair boundary}
The proof on p. 23 infers generic Fourier purity of weight one from generic finite monodromy and the elementary conditions. On p. 25 that purity is used to kill the contribution of a weight-zero term before embedding in the cohomology of a proper curve. Our example exhibits exactly the term that the argument cannot discard.\par
Katz’s cited Theorem 7.3.8(5) was checked in the corrected author manuscript [katz], PDF pp. 211–213 (Chapter 7 local pp. 11–13). Its use of purity is governed by Definition 7.3.7, which requires a middle extension in addition to purity on a lisse open. Definition 7.3.5 likewise distinguishes Fourier sheaves from merely elementary ones. Here \(\mathcal F\) is not a middle extension: its stalk at zero is zero, while \(L_0=\overline{\mathbb Q}_2\). Thus the cited theorem does not supply the purity asserted in the paper.\par
A reasonable repair is to require the appropriate pure middle-extension hypothesis, or to formulate a result for a suitable weight-one quotient. The counterexample and the missing hypothesis are established; a complete corrected realization theorem, including splitting and arbitrary-rank details, is not independently proved in this review. In particular, the counterexample does not invalidate the specific pure middle-extension examples of Katz or the separately cited constructions of Bültel without an additional analysis of those examples.\par
\section*{Two missing restrictions in the density lemma}
Problem 3 (Minor: Lemma 9 needs analyticity and a positive-dimensional density setting).  Lemma 9(1), p. 9, identifies zero density with empty interior for every closed conjugacy-invariant subset of the compact image. This is false for arbitrary closed subsets. In addition, the stated setup on p. 7 allows \(\dim S=0\), where the density conclusions of Lemma 9 fail even for analytic subsets.\par
\subsection*{A positive-dimensional counterexample with an actual density}
Let \(k=\mathbb F_3\), \(S=\mathbb A^1/k\), \(\ell=2\), and let \[\rho:\pi_1(S)\longrightarrow\widehat{\mathbb Z}\longrightarrow\mathbb Z_2
   \longrightarrow\mathrm{GL}_2(\mathbb Q_2),\qquad
 n\longmapsto u(n)=
 \begin{pmatrix}1&n\\0&1\end{pmatrix},\] where the arithmetic projection sends geometric Frobenius to \(1\). Its compact image is isomorphic to \(\mathbb Z_2\), and conjugacy in that image is trivial.\par
Define \[U=\bigcup_{j\geq1}\bigl(-j+2^{j+2}\mathbb Z_2\bigr),
 \qquad C=\mathbb Z_2\setminus U,\qquad \Phi=u(C).\] The negative integers are dense in \(\mathbb Z_2\), so \(U\) is open and dense. Hence \(\Phi\) is closed with empty interior. Its Haar measure satisfies \[\mu(C)\geq1-\sum_{j\geq1}2^{-j-2}=\frac34.\] A Haar-measure calculation alone would not establish the needed Frobenius density. Here it can be established directly. The least positive integer in the \(j\)th removed residue class is \(2^{j+2}-j\). Thus only \(O(\log x)\) classes meet \([1,x]\), and for fixed \(J\) the contribution of the classes with \(j>J\) is at most \[x\sum_{j>J}2^{-j-2}+O(\log x).\] Finite unions of residue classes have their Haar density. Letting first \(x\to\infty\) and then \(J\to\infty\) proves that \(D=C\cap\mathbb Z_{>0}\) has natural density \(d=\mu(C)\geq3/4\).\par
There are \(N_n=3^n/n+O(3^{n/2}/n)\) closed points of degree \(n\) on \(\mathbb A^1_{\mathbb F_3}\), and \(\rho(F_s)=u(n_s)\). With \(\varepsilon=(t-1)\log3\), the numerator in the paper’s density ratio is therefore \[\sum_{n\in D}\frac{e^{-\varepsilon n}}n+O(1),\] whereas the denominator is \(\sum_{n\geq1}e^{-\varepsilon n}/n+O(1)\). Partial summation from the natural density of \(D\) gives ratio tending to \(d\). Consequently \(\delta_S(\Sigma_\Phi)=d\geq3/4\) although \(\Phi\) has empty interior. This is a counterexample under the actual Frobenius density definition, not just an abstract topological objection.\par
\subsection*{The zero-dimensional boundary}
Independently take \(S=\operatorname{Spec}\mathbb F_3\) and the same arithmetic unipotent representation. There is one closed point, of degree one, and the paper’s ratio assigns \(|S|\) density one. Set \(\Phi=\{u(1)\}\). It is a closed analytic subset with empty interior, but \(\Sigma_\Phi=|S|\). Furthermore \(\Phi_{|S|}\) is a singleton, neither dense in \(u(\mathbb Z_2)\) nor Zariski dense in its connected algebraic closure \(\mathbb G_a\). This contradicts parts (1)–(3) of Lemma 9 under the printed zero-dimensional scope.\par
\subsection*{Repair and surviving uses}
For positive-dimensional \(S\), restrict the empty-interior equivalence to closed analytic subsets. Serre’s author lecture manuscript [serre], p. 81, states this hypothesis immediately before Theorem 6.11; p. 80 gives only the general closed-set upper-density bound by Haar measure. This 2011 author text is identified separately from the 2012 edition cited by the paper.\par
The actual subset used in the proof of Lemma 9(2) is explicitly closed analytic, so that positive-dimensional argument survives the restriction. The final assertion in (1), that positive upper density makes the topologically generated subgroup open, can also be obtained from the general Haar bound: the closed generated subgroup has positive Haar measure and hence finite index. It is not refuted by the example above. The zero-dimensional case should be removed from the density machinery and the original local-global questions handled directly when \(S=\operatorname{Spec}k\). These are bounded auxiliary corrections, not a disproof of the main positive-dimensional ghost argument.\par
\section*{Normalizer ambiguity}
Problem 4 (Minor: the conjugating map on zero weights is not canonical).  In the second functoriality bullet of Section 6.4, p. 14, a conjugating element \(g\) with \(\overline T'=g\overline T g^{-1}\) is claimed to be unique up to right multiplication by \(\overline T\). It is then asserted that the resulting map on invariant spaces is independent of \(g\). The ambiguity is the normalizer, not the centralizer.\par
For \(\mathrm{SL}_2\) and its adjoint representation, put \[T=\{\operatorname{diag}(a,a^{-1})\},\quad
 H=\operatorname{diag}(1,-1),\quad
 w=\begin{pmatrix}0&-1\\1&0\end{pmatrix}.\] Then \(V^T=\mathbb Q_\ell H\). Both \(1\) and \(w\) conjugate \(T\) to itself, but they induce \(H\mapsto H\) and \(H\mapsto-H\). Thus the asserted independence fails. This is the rank-one zero-weight situation of the adjoint representation; the large-geometric-monodromy example in Section 3.2, p. 6, produces the corresponding adjoint object when \(g=1\).\par
Choosing a particular conjugating element still induces an isomorphism natural in representations. Accordingly the earlier statement of independence up to natural isomorphism is not refuted. The repair is to remove the uniqueness assertion and record the Weyl-group ambiguity. At the pro-reductive level a proof must also arrange compatible choices in the relevant quotients; that coherence is not supplied merely by the displayed matrix calculation. This finding is local and does not change the ghost’s numerical dimension or its isomorphism class.\par
\section*{Positive checks, dependencies and assessment}
The following parts of the argument were checked at the stated scope.\par
In Section 2, pp. 5–6, an open specialized monodromy image suffices to transport an invariant injective map from a closed fiber to an open subgroup of the generic group. The Hilbertian specialization and Frattini input are invoked dependencies. The descent statement is geometric: the relevant field on p. 6 is \(\overline{k(s)}\), not \(k(s)\). The nonsemisimple Kummer example and the \(\operatorname{End}(V)/\mathbb Q_\ell\operatorname{Id}\) example in Section 3 correctly distinguish actual, almost and semisimplified assertions.\par
Lemma 10, p. 10, uses the convergent logarithm on \(1+\ell M\) (on \(1+4M\) at \(\ell=2\)) to recover a matrix from its positive power. Once that lemma and the corrected density hypothesis are available, Proposition 3’s finite union of closed characteristic-polynomial loci, irreducibility of \(G^\circ\), dominance of the power map at the identity, and semisimple character comparison form a consistent chain.\par
Under the connected and torsion-free reductions in Section 6.2, the finite central kernel in Section 6.3 lies in the geometric torus and acts trivially on its invariants. This explains lift independence, which is distinct from Problem 4’s Weyl ambiguity. In Section 6.5, torsion-freeness permits gcd to commute with taking powers of roots. The degrees of the normalized gcd polynomials decrease and stabilize. The finite set of initial closed points then supplies one common exponent for an almost factor, and Claim 1 yields Theorem 6’s maximality. The split-prime choice, the tracking of twists in Section 6.5.3.4, and the two cases in Claim 2 were examined; the invoked independence of zero-weight dimension, Frobenius-torus theorem, and open-image result were not re-proved.\par
Lemma 13 and Corollary 14, pp. 18–19, do not realize an arbitrary rational pure polynomial as an abelian variety. They take a gcd of actual Frobenius polynomials. If a simple Honda–Tate factor contributes \(\Pi_w^{n_w}\), the gcd exponent is \(n_w\min(m_1,m_2)\) and retains the required divisibility by \(n_w\). After a common constant extension, the normalized polynomials are those of actual base-changed fibers. This mechanism is consistent, conditional on the corrected ghost argument.\par
The implication from a nonzero minuscule weight to the absence of zero weights, and the Rosati argument in Lemma 21, p. 22, were checked. For types I–III the Rosati involution is trivial on the center, so the central subgroup in the polarization kernel satisfies \(h^2=1\) and is finite; this supports the one-dimensional-center and supersingularity conclusion. The characteristic-zero lifting and the specific Bültel/Katz monodromy constructions remain external dependencies. No author code was executed, and no numerical experiment is used as proof.\par
Item | State\par
Overall assessment | Major. Proposition 5’s actual pointwise conclusion and Theorem 22’s stated scope have exact counterexamples.\par
Other findings | Minor restrictions in Lemma 9 and a Minor normalizer ambiguity.\par
Repair scope | Local counterexamples and corrections verified; no complete replacement theorem certified.\par
Central almost theory | Not disproved; imported dependencies and full corrected descent remain outside certification.\par
Scientific decision | Pending; human reviews empty.\par
\section*{Supporting checks}
Primary input: 564150 bytes, SHA-256\par
00000TOKEN\par
The exact file, source images, added primary references and all report-build artifacts are bound in the accompanying inventory. Decisive source pages 2,4,7,9,12–14,23–25 and the cited Katz and Serre pages were actually viewed by the primary executor, together with every final report page. The three native helpers in this conversation supplied bounded text-based checks of the density and early arguments, the central ghost proof, and the abelian/Appendix arguments. Their text reading is not visual inspection or independent acceptance. The primary executor reviewed the mathematical deductions, source-reference hypotheses and report.\par
The supplied manifest and a further bounded title/supplement/correction search identified no dedicated code release or separate erratum. This is a search limitation, not a defect. No publication-version equivalence or precise author revision date is inferred. The Katz corrected author manuscript and Serre’s 2011 author lectures were retrieved separately and are not silently identified with the cited print editions. Source/TeX/compiler/dependency hashes and byte counts before and after the build, real exits, logs, cache changes and actual per-page receipts are retained. Controller UI observations and runtime self-description are recorded separately; backend equivalence is unverified.\par
9 A. Cadoret and A. Tamagawa, Ghosts and families of abelian varieties with a common isogeny factor, author-hosted manuscript, revision date undetermined. https://webusers.imj-prg.fr/\textasciitilde{}anna.cadoret/Ghosts.pdf. Journal DOI: 10.4171/JEMS/1597 (https://doi.org/10.4171/JEMS/1597). N. M. Katz, Exponential sums and differential equations, corrected author manuscript, Chapter 7, local pp. 11–13, Definitions 7.3.4–7.3.7 and Theorem 7.3.8. https://web.math.princeton.edu/\textasciitilde{}nmk/wholebookcorrms.pdf. Exact retrieved PDF pp. 211–213. J.-P. Serre, Lectures on \(N_X(p)\), author lecture manuscript dated 6 September 2011, pp. 80–81, Theorems 6.8 and 6.11. https://www.college-de-france.fr/media/professeurs-honoraires/UPL47720\_cours\_Taiwan\_avril\_\_10juin\_2011\_.pdf. Not a full collation with the cited 2012 print edition.\par
\end{document}
